\chapter{Was Commands wirklich sind}
Das Wort \emph{Command} wird für zwei Dinge benutzt. Erstens gibt es produktseitige Schnellschalter, etwa den dokumentierten Start von Deep Research über \texttt{/Deepresearch}. Zweitens gibt es frei definierte Kurzbefehle wie \texttt{/analyse}. Bei diesen ist der Schrägstrich keine Magie: Er funktioniert, weil seine Bedeutung zuvor als Anweisung vereinbart wurde.
\begin{center}\begin{tikzpicture}[node distance=12mm]
\node[box] (c) {``Command''};\node[tool,below left=of c] (p) {Produktfunktion\\Werkzeug wird gestartet};\node[tool,below right=of c] (k) {Konvention\\Prompt wird expandiert};\draw[arr](c)--(p);\draw[arr](c)--(k);
\end{tikzpicture}\end{center}

\begin{achtung}Verfügbare Menüpunkte, Namen und Limits können sich nach Tarif, Plattform, Region und Rollout unterscheiden. Die aktuelle Oberfläche ist maßgeblich.\end{achtung}
